\section{保持简单：Manus 的产品原则}

前面讨论市场、全球化和基础设施，最后回到最朴素的产品原则。如果把本期压缩成一句话，就是“保持简单”。这不是反功能，而是反稀释。Agent 公司很容易被用户需求、客户场景、大厂竞争和投资人预期拉向复杂；但产品必须有一个用户能记住的核心。

\begin{figure}[H]
\centering
\includegraphics[width=0.96\textwidth]{stay-simple-principle.png}
\caption{保持简单的产品原则：核心价值、克制、特色、增长与长期。}
\end{figure}

\begin{knowledgebox}{读图：简单不是少做，而是守住核心}
图中简单不是停滞，而是围绕核心价值组织增长。功能可以增加，但必须增强而不是稀释特色。Agent 产品尤其需要克制，因为它天然可以接很多工具、做很多任务。
\end{knowledgebox}

\subsection{Manus 的特殊位置}

Manus 最好的预期，是让有高价值工作的白领获得一个 7x24 小时不断推理的工作伙伴。这里的核心不是“更多功能”，而是“持续推理的伙伴”。这和 EP130 张月光讲的 AI 朋友有相通之处：AI 不只是工具，而是一个能持续帮你完成边界外工作的对象。

\begin{importantbox}{从工具到伙伴}
Agent 如果只是按钮集合，就会被功能列表拖累；如果成为持续推理的伙伴，它的价值来自关系、信任和稳定产出。
\end{importantbox}

\subsection{本章小结}

保持简单，是 Manus 增长后的核心挑战。越强的 Agent 越容易复杂，越需要守住用户能记住的核心价值。

